
\section{Evaluation Metrics}
\label{sec:evaluation_metrics}

Telemetry data is collected from the workload generator, Sequencer, Executor, Prover, and Submitter during each benchmark run. These components emit raw JSON, JSONL, and CSV files into the \path{benchmark-suite/metrics/} directory. The Data Tools analytics pipeline then processes these outputs using Python aggregation scripts to compute run-level metrics for throughput, latency, finalization, proving overhead, data availability footprint, and cost.

The evaluation metrics are selected to capture both general distributed-system performance and rollup-specific trade-offs. Throughput and latency measure the user-facing performance of the Layer-2 pipeline. Finalization time measures how long a sealed batch takes to reach Layer-1 settlement. Gas and cost metrics capture the economic impact of submitting batches to Layer 1. Data availability metrics measure the payload footprint that directly affects calldata or blob-related cost. Table~\ref{tab:metrics} summarizes the main metrics used in the evaluation.

\begin{table}[!htbp]
\centering
\caption{Performance Metrics Used in Benchmarking}
\label{tab:performance_metrics}
\footnotesize
\setlength{\tabcolsep}{4pt}
\renewcommand{\arraystretch}{1.18}
\begin{tabularx}{\textwidth}{@{}p{0.26\textwidth}p{0.14\textwidth}Y@{}}
\toprule
\textbf{Metric} & \textbf{Unit} & \textbf{Definition and purpose} \\
\midrule

Throughput 
& TPS 
& Measures the rate of successfully processed transactions during a benchmark run. It indicates the effective processing capacity of the rollup pipeline. \\

\addlinespace

Run success rate 
& \% 
& Measures the fraction of submitted transactions that are accepted and successfully processed. It helps identify whether the system remains stable under the offered workload. \\

\addlinespace

Soft finality latency 
& ms 
& Measures the time from transaction submission or Sequencer intake until the transaction is included in a sealed batch. It represents the first user-visible confirmation point. \\

\addlinespace

Hard finality latency 
& ms 
& Measures the time from batch sealing or submission until the corresponding update is settled on Layer 1. It captures the delay between Layer-2 batching and Layer-1 finalization. \\

\addlinespace

Proof time 
& ms / s 
& Measures the time spent in mock proof delay or real prover execution. It is used to identify whether the Prover becomes the main bottleneck. \\

\bottomrule
\end{tabularx}
\normalsize
\end{table}

\begin{table}[!htbp]
\centering
\caption{Cost and Data Availability Metrics Used in Benchmarking}
\label{tab:cost_da_metrics}
\footnotesize
\setlength{\tabcolsep}{4pt}
\renewcommand{\arraystretch}{1.18}
\begin{tabularx}{\textwidth}{@{}p{0.26\textwidth}p{0.14\textwidth}Y@{}}
\toprule
\textbf{Metric} & \textbf{Unit} & \textbf{Definition and purpose} \\
\midrule

Layer-1 gas used 
& Gas units 
& Measures the gas consumed by the Layer-1 batch submission transaction. It is the most stable metric for comparing settlement overhead across configurations. \\

\addlinespace

Cost per tx 
& Wei / USD 
& Estimates the total batch cost divided by the number of transactions in the batch. Wei-based values are used for stable comparison, while USD values are treated as explanatory estimates. \\

\addlinespace

DA footprint 
& Bytes 
& Measures the size of the encoded or compressed batch payload submitted through the selected data availability mode. It directly affects calldata or blob-related cost. \\

\addlinespace

Compression ratio 
& Ratio 
& Measures the reduction between raw transaction data and the compressed or encoded payload. It indicates the effectiveness of payload reduction before data availability submission. \\

\addlinespace

Blob utilization 
& \% 
& Measures how much of the target blob capacity is filled by the batch payload. It is mainly used when evaluating EIP-4844-style blob submission. \\

\bottomrule
\end{tabularx}
\normalsize
\end{table}


In addition to average values, the analysis reports distributional statistics because latency and queueing behaviour in distributed systems are often skewed. The reported statistics include mean, median, standard deviation, 95\% confidence intervals, and latency percentiles such as P50, P90, P95, and P99. These values are particularly important for identifying tail-latency effects under high-load and burst-load conditions.

The aggregation pipeline also links metrics across different telemetry sources. For example, Sequencer metrics describe queueing and batch formation, Executor metrics describe execution and proving time, and Submitter metrics describe Layer-1 gas usage, data availability footprint, and submission latency. Combining these metrics allows the evaluation to identify whether a bottleneck originates from batching, execution, proving, data availability formatting, or Layer-1 submission.
